%% --- Appendix F: Discovery Scaling ---

\section{Solution Discovery: Search Scaling}
\label{sec:app:discovery-scaling}

Table~\ref{tab:solver} reports \sys{} scores from the final iteration of the best-performing branch (1 of 5 parallel branches).
Table~\ref{tab:scaling} reports the best score across \emph{all} nodes in the search tree, measuring the discovery module's ceiling under varying tree and budget configurations. Solutions are manually checked and specification-violating ones (metric gaming on LLM-SQL, abusing the Prism scoring formula by failing on subsets of test inputs) are excluded. All scaling configurations report a single seed and cross-seed variance can be substantial (e.g., TXN ranges from 3636 to 3906 across three seeds with the base configuration). So \textbf{these results should be interpreted as directional rather than definitive}.

% \sys{}'s discovery module uses a parallel explore-exploit search tree parameterized by depth (iterations $I$), width (parallel branches $B$), and pruning (branches retained per iteration, $K$). The base configuration ($I\!=\!5, B\!=\!5, K\!=\!2$) retains the top-2 branches at each iteration, with a per-node evaluator budget of $E\!=\!4$ (the agent may attempt up to 4 evaluated solution versions per node, subject to early stopping via patience).

\paragraph{Tree and budget scaling.}
We evaluate search configurations along three axes: width (parallel branches $B$), depth (iterations $I$), and per-node evaluator budget ($E$, maximum evaluated solution versions per node).
Table~\ref{tab:scaling} reports the best score across all nodes for each configuration (single run each).
The tree scaling block in Table~\ref{tab:scaling} varies the tree structure ($I$, $B$, $K$) at fixed budget ($E\!=\!4$). The budget scaling block varies $E$ while holding the tree fixed at $I\!=\!5, B\!=\!5, K\!=\!2$.

\begin{table}[h]
\centering
\small
\caption{Best score across all search tree nodes under different tree and budget configurations (single run each). $I$=iterations (depth), $B$=branches (width), $K$=branches retained per iteration, $E$=max evaluator calls per node. ``Base'' is the configuration used for Table~\ref{tab:solver}. \textbf{Bold} marks the best score in each column, computed separately for baselines and \sys{} configurations.}
% $^\ddagger$Reference baselines are \emph{selector-based}: each system's own pipeline chose which solution to use for paper writing (e.g., Sakana selects via its agent, \sys{} selects the final iteration of the best-performing branch). \sys{} scaling configurations report the best score across all nodes, independent of the selection mechanism.
\label{tab:scaling}
\setlength{\tabcolsep}{4pt}
\begin{tabular}{@{}lcccccccccc@{}}
\toprule
\textbf{Config} & $I$ & $B$ & $K$ & $E$ & \textbf{Nodes} & \textbf{Prism}$\uparrow$ & \textbf{Cloud.}$\downarrow$ & \textbf{EPLB}$\uparrow$ & \textbf{LLM-SQL}$\uparrow$ & \textbf{TXN}$\uparrow$ \\
\midrule
\multicolumn{11}{@{}l}{\textit{Reference baselines (best-of-3 seeds, selector-based; from Table~\ref{tab:solver})}} \\
Human           & \multicolumn{5}{c}{--} & 21.89 & 626.24 & 0.1265 & 0.6920 & 2724.8 \\
AdaEvolve       & \multicolumn{5}{c}{--} & \textbf{26.26} & 637.10 & 0.1450 & \textbf{0.7520} & 4310 \\
EvoX            & \multicolumn{5}{c}{--} & \textbf{26.26} & 623.69 & \textbf{0.1453} & 0.7300 & 4310 \\
Sakana          & \multicolumn{5}{c}{--} & \textbf{26.26} & 627.11 & 0.1270 & 0.7320 & 4184 \\
ARC             & \multicolumn{5}{c}{--} & 26.25 & 690.37 & 0.1266 & 0.6757 & 3247 \\
AIR             & \multicolumn{5}{c}{--} & \textbf{26.26} & 734.28 & 0.1449 & 0.7148 & \textbf{4311} \\
DS              & \multicolumn{5}{c}{--} & \textbf{26.26} & \textbf{620.09} & 0.1284 & 0.7307 & 4286 \\
\midrule
\multicolumn{11}{@{}l}{\textit{\sys{} scaling configurations (1st seed of 3, best across all nodes)}} \\
Base            & 5 & 5 & 2  & 4  & 25  & 26.26 & 618.09 & 0.1287 & 0.7299 & 3636 \\
\midrule
\multicolumn{11}{@{}l}{\textit{Tree scaling (fixed budget $E\!=\!4$)}} \\
No pruning      & 5 & 5 & 5  & 4  & 25  & 26.26 & 618.08 & \textbf{0.1461} & 0.7092 & 3663 \\
Wide            & 5 & 10 & 2 & 4  & 50  & 26.36 & 618.09 & 0.1369 & 0.7215 & 4082 \\
Wider           & 5 & 15 & 2 & 4  & 75  & 26.26 & 618.08 & 0.1456 & 0.7189 & 4237 \\
Widest          & 5 & 20 & 2 & 4  & 100 & \textbf{26.44} & \textbf{618.07} & 0.1455 & 0.7257 & 4255 \\
Deep            & 10 & 5 & 2 & 4  & 50  & \textbf{26.44} & 618.08 & 0.1460 & 0.7118 & 3831 \\
Deep+wide       & 10 & 10 & 4 & 4  & 100 & 26.33 & 618.09 & 0.1458 & 0.7078 & 4000 \\
\midrule
\multicolumn{11}{@{}l}{\textit{Budget scaling (fixed tree $I\!=\!5, B\!=\!5, K\!=\!2$)}} \\
Budget 200      & 5 & 5 & 2  & 8  & 25  & 26.26 & \textbf{618.07} & 0.1284 & 0.7256 & \textbf{4348} \\
Budget 500      & 5 & 5 & 2  & 20 & 25  & 26.34 & \textbf{618.07} & 0.1456 & \textbf{0.7316} & 4237 \\
% \midrule
% \multicolumn{11}{@{}l}{\textit{Combined}} \\
% Wide+2$\times$budget & 5 & 10 & 4 & 8  & 50  & 26.50 & 618.08 & 0.1462 & 0.7276 & \textbf{3831} \\
\bottomrule
\end{tabular}
\end{table}

Three patterns emerge from tree scaling.
First, \textbf{TXN scales monotonically with width}: scores increase steadily from 3636 (base, $B\!=\!5$) through 4082 ($B\!=\!10$), 4237 ($B\!=\!15$), and 4255 ($B\!=\!20$), a 17\% improvement at the widest configuration and approaching AdaEvolve (4310).
Second, \textbf{EPLB benefits from scale but saturates early}: most non-base configurations reach ${\sim}0.146$ (a 13\% improvement over the base's 0.129), with the exception of Wide ($B\!=\!10$, $K\!=\!2$) at 0.137.
Third, \textbf{Cloudcast, LLM-SQL, and Prism largely saturate}: all configurations converge to similar scores on these three tasks regardless of tree shape, suggesting a narrow basin of high-performing solutions that the default search finds quickly.

%\paragraph{Budget scaling.}
%Budget scaling reveals a split between tasks.
%\textbf{TXN responds strongly to increased budget}: budget~200 reaches 4348 (+20\% over the base), comparable to the widest tree configuration ($B\!=\!20$, 4255).
%Interestingly, budget~500 does not improve further (4237), suggesting that TXN gains saturate around $E\!=\!8$ per node and that the additional budget is spent refining rather than discovering better algorithmic strategies.
%EPLB shows a similar pattern: budget~500 reaches 0.1456 (+13\%), matching the width-scaled results but with only 25 nodes instead of 75--100.
%Cloudcast, LLM-SQL, and Prism remain flat across all budget levels, consistent with the tree scaling results.

Taken together, \textbf{width (more independent branches) is the most efficient scaling axis} for tasks with diverse solution strategies.
The widest tree ($B\!=\!20$, 100~nodes, $E\!=\!4$) matches or exceeds the highest per-node budget configuration ($E\!=\!20$, 25~nodes) on 4 of 5 tasks, despite using 5$\times$ fewer evaluator calls per node.
Budget scaling shows gains on TXN (budget~200 reaches 4348, +20\%) but saturates quickly. Budget~500 does not improve further on TXN, and the remaining tasks are flat across all budget levels.
Depth and budget produce diminishing returns once the search covers the main algorithmic strategies.

%\paragraph{Comparability with AdaEvolve's iteration budget.}
%AdaEvolve reports results under a fixed iteration budget of $T\!=\!100$~\citep{cemri2026adaevolve}.
%Our base configuration produces ${\sim}100$ evaluator calls (25 nodes $\times$ ${\sim}4$ calls each), making the evaluator-call counts directly comparable; the difference is that \sys{} distributes evaluations across a tree with ideation and pruning between iterations, adding LLM overhead not counted in the evaluator budget.

\paragraph{Specification violations at higher budgets.}
Increasing per-node budget amplifies specification violation risk.
At budget~100, no specification violations are observed on Prism. At budget~200 and~500, 2--8\% of nodes converge to solutions that exploit the scoring formula rather than solving the task correctly.
LLM-SQL shows a similar trend: the fraction of nodes flagged for metric gaming by the post-hoc auditor grows from ${\sim}0\%$ at budget~100 to ${\sim}50\%$ at budget~200 and ${\sim}70\%$ at budget~500.
In contrast, wider trees at budget~100 show lower violation rates despite producing more total nodes, because each node has fewer iterations to discover and refine exploitative patterns.


%Scores in Table~\ref{tab:scaling} reflect the best evaluated version at each node.
%Within a single node, the agent iterates through multiple solution versions, and gold-evaluator verification reveals that earlier versions sometimes outperform the final one---for example, one node in the no-pruning configuration reaches Prism~31.73 at version~2 before regressing to 26.26 at version~4.
%This version-level selection loss (the agent iterating past its best solution) is distinct from the node-level selection reported here and represents an orthogonal improvement opportunity in the discovery module's within-node search strategy.

